\documentclass[11pt]{article}

\usepackage[T1]{fontenc}
\usepackage{lmodern}
\usepackage{amsmath,amssymb,mathtools}
\usepackage{booktabs}
\usepackage{enumitem}
\usepackage[hidelinks]{hyperref}
\usepackage[margin=1in]{geometry}
\usepackage{xcolor}

\newcommand{\Hilbert}{\mathcal H}
\newcommand{\C}{\mathbb C}
\newcommand{\ii}{\mathrm i}
\newcommand{\hc}{\mathrm{h.c.}}
\newcommand{\ket}[1]{\lvert #1\rangle}
\newcommand{\bra}[1]{\langle #1\rvert}
\newcommand{\braket}[2]{\langle #1\vert #2\rangle}
\newcommand{\projector}[1]{\lvert #1\rangle\!\langle #1\rvert}

\title{From Bloch States to Wannier Functions\\
\large Gauge, Localization, and Operator Reconstruction}
\author{PhysKit exploratory lecture notes}
\date{}

\begin{document}
\maketitle

\begin{center}
\fcolorbox{black}{gray!10}{\parbox{0.9\linewidth}{
\textbf{Artifact status.} These notes accompany the accepted exploratory notebook
\texttt{bloch2wannier.ipynb}. They are not a canonical or
supported PhysKit artifact and make no physical- or pedagogical-validation claim.
}}
\end{center}

\section{Scope and learning objectives}

We study a finite one-dimensional, two-orbital, periodic tight-binding model.
The construction separates the physical operator, its matrix representations,
the retained band subspace, the gauge-dependent Bloch frame, and the resulting
Wannier basis.

After working through the notebook and these notes, a reader should be able to:
\begin{enumerate}[label=\arabic*.]
  \item construct the finite periodic Hamiltonian used in the notebook;
  \item decompose it into Bloch fibers with a discrete Fourier transform;
  \item distinguish Bloch eigenvectors from band projectors;
  \item explain how momentum-dependent phases change the Wannier basis;
  \item construct and check translated Wannier functions;
  \item reconstruct the retained dispersion from Wannier matrix elements; and
  \item state why independently constructed retained spaces require alignment.
\end{enumerate}

\section{Finite periodic state space}

For $N$ unit cells with orbitals $A$ and $B$ in every cell,
\begin{equation}
  \Hilbert_{\mathrm{real}}=\C^N\otimes\C^2,
  \qquad \dim\Hilbert_{\mathrm{real}}=2N.
\end{equation}
The orthonormal basis is
$\{\ket{R,A},\ket{R,B}\}_{R=0}^{N-1}$, with cell labels understood modulo $N$.
This is a finite representation of a periodic model, not a material-orbital
geometry claim.

\section{Real-space Hamiltonian}

The notebook uses the SSH-like Hamiltonian
\begin{align}
\hat H={}&\sum_R\Delta\bigl(\projector{R,A}-\projector{R,B}\bigr)\notag\\
&+\sum_R\left[
 t_1\ket{R,A}\bra{R,B}
 +t_2\ket{R+1,A}\bra{R,B}+\hc\right].
\end{align}
Here $t_1$ and $t_2$ are intracell and intercell hopping amplitudes, and
$\pm\Delta$ are onsite energies. The notebook assembles the $2N\times2N$ matrix
and verifies Hermiticity numerically. That is an internal consistency check, not
physical validation of the model for a material.

\section{Translation symmetry and Bloch fibers}

The allowed momenta and Fourier matrix are
\begin{equation}
 k_j=\frac{2\pi j}{Na},
 \qquad
 F_{kR}=\frac{1}{\sqrt N}e^{-\ii kRa}.
\end{equation}
With $\mathcal F=F\otimes I_2$,
\begin{equation}
 H_k^{\mathrm{full}}=\mathcal F H_{\mathrm{real}}\mathcal F^\dagger.
\end{equation}
Translation symmetry makes this matrix block diagonal. Under the notebook's
conventions, each $2\times2$ block is
\begin{equation}
 H(k)=
 \begin{pmatrix}
  \Delta & t_1+t_2e^{-\ii ka}\\
  t_1+t_2e^{\ii ka} & -\Delta
 \end{pmatrix}.
\end{equation}
The executable notebook checks Fourier unitarity, equality of every transformed
block with $H(k)$, and vanishing of all off-fiber blocks.

\section{Eigenstates, projectors, and gauge freedom}

At each momentum,
\begin{equation}
 H(k)\ket{u_{nk}}=E_n(k)\ket{u_{nk}}.
\end{equation}
Retaining the lower band gives a normalized vector $\ket{u_k}$ and rank-one
projector
\begin{equation}
 P(k)=\projector{u_k}.
\end{equation}
The vector is a choice of frame inside the retained subspace. For an isolated
band,
\begin{equation}
 \ket{u_k}\longmapsto e^{\ii\phi(k)}\ket{u_k}
\end{equation}
changes that frame while leaving $P(k)$ unchanged. The notebook demonstrates
this with a seeded irregular phase field and an executable projector check.

\section{Open-chain parallel transport}

A smoother sampled frame is built sequentially. The phase of
$\ket{u_{k_{j+1}}}$ is chosen so that
\begin{equation}
 \braket{u_{k_j}}{u_{k_{j+1}}}
\end{equation}
is real and positive. Neighboring overlap phases are thereby removed along the
ordered samples.

\paragraph{Boundary of the construction.}
This is open-chain parallel transport. It does not enforce closure between the
last and first momentum samples. The remaining closure phase must be addressed
explicitly before making a periodic-gauge or Berry-phase claim.

\section{Wannier construction}

With $\ket{\psi_k}=\ket{k}\otimes\ket{u_k}$, define
\begin{equation}
 \ket{w_{R_0}}=\frac{1}{\sqrt N}\sum_k
 e^{-\ii kR_0a}\ket{\psi_k}.
\end{equation}
Its real-space orbital amplitudes are
\begin{equation}
 \braket{R,\alpha}{w_{R_0}}
 =\frac1N\sum_k e^{\ii k(R-R_0)a}u_\alpha(k),
\end{equation}
and its cell probability is
\begin{equation}
 p_{R_0}(R)=\sum_\alpha
 \left|\braket{R,\alpha}{w_{R_0}}\right|^2.
\end{equation}
Different phase frames produce different shapes for individual Wannier
functions while spanning the same retained subspace.

\section{Wannier-basis consistency}

Place the translated Wannier functions in the columns of a $2N\times N$ matrix
$W$. The notebook checks:
\begin{align}
 W^\dagger W &= I_N &&\text{(orthonormality)},\\
 WW^\dagger &= P_{\mathrm{Bloch}} &&\text{(retained-projector equality)},
\end{align}
as well as translation covariance between adjacent center labels. These checks
establish consistency of the finite representations, not completeness outside
the retained band.

\section{Exploratory localization diagnostic}

Using minimum-image displacements $x_R$ relative to $R_0$, the notebook computes
\begin{equation}
 \Omega_{R_0}=\sum_Rp_{R_0}(R)x_R^2
 -\left(\sum_Rp_{R_0}(R)x_R\right)^2.
\end{equation}
This diagnostic permits a finite-example comparison between random and
parallel-transport frames.

\paragraph{Boundary of the diagnostic.}
$\Omega_{R_0}$ is not a periodic-position functional and is not the
Marzari--Vanderbilt spread. A smaller value in this example is not a
maximal-localization result.

\section{Hamiltonian in Wannier coordinates}

The retained Hamiltonian is
\begin{equation}
 H_{RR'}=\bra{w_R}\hat H\ket{w_{R'}},
 \qquad H_W=W^\dagger H_{\mathrm{real}}W.
\end{equation}
Translation symmetry makes $H_W$ circulant. A row referenced to an origin cell
defines $H(\Delta R)$, from which the retained dispersion is reconstructed:
\begin{equation}
 E(k)=\sum_{\Delta R}e^{\ii k\Delta Ra}H(\Delta R).
\end{equation}
The notebook checks this reconstruction against the lower-band eigenvalues at
every sampled momentum.

\section{The special single-band result}

For one isolated band,
\begin{equation}
 H_{RR'}=\frac1N\sum_k e^{\ii k(R-R')a}E(k),
\end{equation}
so the $U(1)$ phase does not appear. Let
$S=W_{\mathrm{smooth}}^\dagger W_{\mathrm{random}}$. The basis change can be
nontrivial while
\begin{equation}
 H_{\mathrm{random}}=S^\dagger H_{\mathrm{smooth}}S,
 \qquad [H_{\mathrm{smooth}},S]=0
\end{equation}
within the notebook's numerical tolerance. Consequently, the two single-band
Wannier Hamiltonians are elementwise equal in this experiment.

\paragraph{Boundary of the result.}
This single-band $U(1)$ experiment cannot demonstrate a nonzero
Wannier-Hamiltonian coordinate difference. A composite retained subspace with
momentum-dependent $U(M)$ mixing would be a different construction and remains
deferred.

\section{Motivation for retained-space alignment}

If pristine and doped calculations independently produce
$\hat H_b^{(P)}$ and $\hat H_d^{(P)}$, their matrices are not automatically in
identified coordinates. An explicit unitary map
\begin{equation}
 \hat U_d:\Hilbert_b^{(P)}\longrightarrow\Hilbert_d^{(P)}
\end{equation}
is required before defining
\begin{equation}
 \Delta\hat H_d^{(P)}
 =\hat U_d^\dagger\hat H_d^{(P)}\hat U_d-\hat H_b^{(P)}.
\end{equation}
The notebook motivates this expression but does not construct two distinct
retained spaces, produce $\hat U_d$, or validate a pristine--dopant subtraction.

\section{Representation map and lecture sequence}

The construction can be summarized as
\begin{equation}
\Hilbert_{\mathrm{real}}
\xrightarrow{\mathcal F}
\bigoplus_k\Hilbert(k)
\xrightarrow{P(k)}
\bigoplus_k\Hilbert^{(P)}(k)
\xrightarrow{\mathrm{gauge}}
\{\ket{u_k}\}
\xrightarrow{\mathrm{Fourier}}
\{\ket{w_R}\}.
\end{equation}

A suggested lecture order is:
\begin{enumerate}
 \item build the real-space basis and matrix;
 \item derive and verify the Bloch fibers;
 \item contrast eigenvectors with projectors;
 \item demonstrate gauge freedom and open-chain smoothing;
 \item construct and compare Wannier functions;
 \item check the basis and reconstruct the band;
 \item explain the special single-band invariance; and
 \item motivate, but do not claim, retained-space alignment.
\end{enumerate}

\section{Preserved limitations}

\begin{enumerate}
 \item The open-chain parallel-transport gauge does not enforce periodic
 momentum closure.
 \item The minimum-image quadratic spread is exploratory, not a
 periodic-position or Marzari--Vanderbilt localization functional.
 \item The single-band $U(1)$ experiment cannot demonstrate a nonzero
 Wannier-Hamiltonian coordinate difference; multiband redesign is deferred.
 \item The pristine--dopant discussion is motivation only and does not construct
 or validate retained-space alignment.
\end{enumerate}

Periodic-gauge closure, Zak phase, a robust periodic spread functional,
composite-band $U(M)$ transformations, entangled-band selection,
Marzari--Vanderbilt localization, imported electronic-structure data, comparison
with Wannier90, and formal PhysKit capability status all require separate
authorization.

\end{document}
